\section{Conclusion}
\label{sec:conclusion}

We presented \method, a standardized protocol for evaluating NAS methods under distribution shift.
Our experiments on the full NAS-Bench-201 search space reveal a striking finding: while Spearman rank correlations across the architecture space appear reassuringly high ($\rho > 0.7$), the top-100 architectures by clean accuracy overlap only 34\% with cross-dataset rankings and 10\% with cross-domain rankings---dropping to 0\% at $K{=}10$.
Top-weighted Kendall-$\tau$ is \emph{negative} ($\tau = -0.41$), demonstrating that among architectures that matter, in-distribution and cross-domain performance are inversely correlated.

The proposed SASC-Pool addresses this with a practical, pool-based criterion that uses z-score normalization over the observed candidate pool rather than requiring access to the full search space.
End-to-end evaluation on four NAS algorithms shows that SASC-Pool consistently improves cross-domain generalization by +0.4--2.8\% on ImageNet-16-120 at zero additional evaluation cost.
Real CIFAR-10-C evaluation on 46 architectures confirms these findings: top architectures lose up to 51\% accuracy under corruption, and ranking by clean accuracy differs from ranking by corrupted accuracy (60\% top-10 overlap).

\paragraph{Limitations.}
Our experiments use the NAS-Bench-201 search space, which is relatively small (15,625 architectures) compared to modern NAS spaces.
The NAS algorithm simulations use proxy scores rather than running actual weight-sharing or gradient-based search.
Our real CIFAR-10-C validation covers 46 of the 15,625 architectures; extending to the full space would enable direct SASC comparison on real corruption data.
Our protocol currently focuses on image classification; extending to detection, segmentation, and other tasks is important future work.

\paragraph{Future Work.}
We plan to extend \method to larger search spaces (NAS-Bench-301, DARTS space), scale the real CIFAR-10-C evaluation to cover more architectures, and develop \sasc variants with adaptive weight selection based on the deployment scenario.
We also aim to release pre-computed corruption metrics as an extension to existing NAS benchmarks, lowering the barrier to shift-aware NAS evaluation.
